\section{Objecthood under strict contraction}\label{sec:objecthood}

This section proves Theorem~\ref{thm:NG_OBJECT_CONTRACTIVE}.
The theorem is formulated on a finite simplex equipped with total variation distance and a strict Dobrushin contraction bound.
Within this admissible regime, the induced Markov operator has a unique fixed distribution, its iterates converge geometrically, and any two $\varepsilon$-approximate fixed points must lie within an explicit separation bound.

Let $Y$ be a finite set.
For $\mu,\nu\in\Delta(Y)$, define the total variation distance by
\[
d_{\mathrm{TV}}(\mu,\nu) := \frac12\sum_{y\in Y}|\mu(y)-\nu(y)|.
\]
We also use the equivalent variational form
\[
d_{\mathrm{TV}}(\mu,\nu)=\sup_{A\subseteq Y}(\mu(A)-\nu(A)).
\]
Let $P:Y\to\Delta(Y)$ be a Markov kernel and let $T:\Delta(Y)\to\Delta(Y)$ be the induced affine map
\[
T(\mu):=\mu P.
\]
The Dobrushin coefficient of $P$ is
\[
\lambda(P) := \max_{y,y'\in Y} d_{\mathrm{TV}}\bigl(P(y,\cdot),P(y',\cdot)\bigr).
\]
This contraction coefficient and its standard consequences for finite chains are discussed in \citep{Dobrushin1956,Seneta2006}.

\begin{lemma}[Dobrushin contraction in total variation]\label{lem:objecthood_dobrushin}
For all $\mu,\nu\in\Delta(Y)$,
\[
d_{\mathrm{TV}}(\mu P,\nu P)\le \lambda(P)\,d_{\mathrm{TV}}(\mu,\nu).
\]
\end{lemma}

\begin{proof}
See Appendix~\ref{app:probkl}.
\end{proof}

\begin{lemma}[Strict contraction gives a unique fixed distribution]\label{lem:objecthood_unique_fixed}
Assume $\lambda(P)<1$.
Then there exists a unique distribution $\pi\in\Delta(Y)$ with $\pi P=\pi$.
Moreover, for every $\mu\in\Delta(Y)$ and every integer $t\ge 0$,
\[
d_{\mathrm{TV}}(\mu P^t,\pi)\le \lambda(P)^t\,d_{\mathrm{TV}}(\mu,\pi).
\]
\end{lemma}

\begin{proof}
See Appendix~\ref{app:probkl}.
\end{proof}

\begin{lemma}[Approximate fixed-point separation bound]\label{lem:objecthood_approximate_separation}
Assume $\lambda(P)\le \lambda<1$.
If $\mu,\nu\in\Delta(Y)$ satisfy
\[
d_{\mathrm{TV}}(\mu P,\mu)\le \varepsilon,
\qquad
d_{\mathrm{TV}}(\nu P,\nu)\le \varepsilon,
\]
then
\[
d_{\mathrm{TV}}(\mu,\nu)\le \frac{2\varepsilon}{1-\lambda}.
\]
\end{lemma}

\begin{proof}
See Appendix~\ref{app:probkl}.
\end{proof}

\begin{proof}[Proof of Theorem~\ref{thm:NG_OBJECT_CONTRACTIVE}]
The uniqueness and convergence claims follow from Lemma~\ref{lem:objecthood_unique_fixed} once $\lambda(P)<1$.
For the approximate objecthood statement, suppose $\mu$ and $\nu$ are $\varepsilon$-stable packaged states in the admissible TV/Dobrushin regime.
Lemma~\ref{lem:objecthood_approximate_separation} then gives the bound
\[
d_{\mathrm{TV}}(\mu,\nu)\le \frac{2\varepsilon}{1-\lambda(P)}.
\]
In particular, when $\varepsilon=0$ every fixed distribution coincides with the unique stationary distribution from Lemma~\ref{lem:objecthood_unique_fixed}.
Thus the theorem rules out multiplicity of well-separated robust objects inside the stated contraction regime.
\end{proof}

\begin{remark}[Regime and metric scope]
This front is formulated in total variation distance with the Dobrushin coefficient as the contraction parameter.
The theorem asserts an obstruction only in that regime.
Other metrics, other contraction notions, or noncontractive operators lie outside the scope of this statement.
\end{remark}

\begin{remark}[Formalization status]
For this front, the formalized component is the uniqueness consequence of strict contraction for fixed distributions.
The total-variation/Dobrushin framework and the $\varepsilon$-stability separation bound are used here at the level of the paper proof and are not part of the formalized core.
\end{remark}
